\documentclass[11pt,a4paper]{article}
\usepackage[T1]{fontenc}
\usepackage[utf8]{inputenc}
\usepackage{amsmath,amssymb,amsthm}
\usepackage[margin=1in]{geometry}
\usepackage[colorlinks=true,linkcolor=blue,urlcolor=blue]{hyperref}
\theoremstyle{plain}
\newtheorem{theorem}{Theorem}
\newtheorem{lemma}[theorem]{Lemma}
\newtheorem{proposition}[theorem]{Proposition}
\newtheorem{corollary}[theorem]{Corollary}
\theoremstyle{definition}
\newtheorem{remark}[theorem]{Remark}

\title{The half-window of the powers of 3:\\
a proof of Karttunen's conjecture on OEIS A037097,\\
in the corrected form}
\author{Adrian Perez Fontelles\\ \small Independent researcher}
\date{24 August 2026}
\begin{document}
\maketitle

\begin{abstract}
OEIS A037097 reads one period of the $n$-th binary column of the powers of 3 --- a window
of length $2^{n-1}$ --- as a polynomial over $\mathrm{GF}(2)$. A.~Karttunen conjectured in
January 1999 that this polynomial is divisible by a large power of $x+1$. The exponent as
originally posted was $2^{n-1}-1$; it was corrected to $2^{n-2}-1$ on the entry by Q.~Gu on
24 July 2026. We prove the corrected statement, and show the exponent is exact:
\[
  a(n)=(x+1)^{2^{n-2}-1}\bigl(B(x)(x+1)+x^{2^{n-2}}\bigr),\qquad n\ge3,
\]
where $B$ is the first quarter of the doubled window. We also confirm that the
pre-correction exponent was indeed too large, already at $n=3$. The proof rests on one
2-adic fact: $3^{2^{n-2}}\equiv1+2^{n}\pmod{2^{n+1}}$, so advancing the exponent by
$2^{n-2}$ complements bit $n$.
\end{abstract}

\section{The sequence and the conjecture}

OEIS A037097 (Antti Karttunen, 29 January 1999), \emph{Periodic vertical binary vectors of
powers of 3, starting from bit-column 2 (halved)}, has offset 2 and
\[
  a(n)=\sum_{k=0}^{2^{n-1}-1}\Bigl(\bigl\lfloor 3^{k}/2^{n}\bigr\rfloor \bmod 2\Bigr)2^{k},
  \qquad a(2),\dots,a(5)=0,\;12,\;120,\;57120.
\]
Identifying an integer with the $\mathrm{GF}(2)$ polynomial whose coefficient of $x^{k}$ is
bit $k$, the entry carries

\begin{quote}
\textbf{Conjecture:} For $n \ge 3$, each term $a(n)$, when considered as a
$\mathrm{GF}(2)[X]$-polynomial, is divisible by the $\mathrm{GF}(2)[X]$-polynomial
$(x+1)^{\mathrm{A000225}(n-2)}$ ($=\mathrm{A051179}(n-2)$). [corrected by Qiyuan Gu,
Jul 24 2026]
\end{quote}

Here $\mathrm{A000225}(n-2)=2^{n-2}-1$; the parenthetical is consistent, since
$(x+1)^{2^{m}-1}=1+x+\dots+x^{2^{m}-1}$ encodes as the integer $2^{2^{m}}-1
=\mathrm{A051179}(m)$. As of the ``Last modified'' line on the live entry (23 August 2026)
the statement is still recorded as an unproven conjecture, with no proof in the entry's
links or programs.

Throughout $n\ge3$ is fixed,
\[
  q=2^{n-2},\qquad b(k)=\bigl\lfloor 3^{k}/2^{n}\bigr\rfloor\bmod 2,\qquad
  a(n)\longleftrightarrow A(x)=\sum_{k=0}^{2q-1}b(k)x^{k},
\]
so the window has length $2^{n-1}=2q$. For nonzero $P\in\mathrm{GF}(2)[x]$ let $v(P)$
denote the exponent of $x+1$ in $P$. We use $x^{2^{m}}+1=(x+1)^{2^{m}}$, the Frobenius
endomorphism in characteristic 2.

\section{The flip}

\begin{lemma}\label{lem:lte}
For $m\ge1$, $v_{2}\bigl(3^{2^{m}}-1\bigr)=m+2$; hence
$3^{2^{m}}\equiv1+2^{m+2}\pmod{2^{m+3}}$, and the multiplicative order of $3$ modulo
$2^{m+3}$ is $2^{m+1}$.
\end{lemma}

\begin{proof}
Lifting the exponent at $p=2$ with even exponent $N$ gives
$v_{2}(3^{N}-1)=v_{2}(3-1)+v_{2}(3+1)+v_{2}(N)-1=1+2+v_{2}(N)-1$; take $N=2^{m}$. Writing
$3^{2^{m}}-1=2^{m+2}u$ with $u$ odd gives the congruence, and the order statement follows
since $3^{2^{m}}\not\equiv1$ but $3^{2^{m+1}}\equiv1$ modulo $2^{m+3}$.
\end{proof}

With $m=n-2\ge1$, read modulo $2^{n+1}$: the order of 3 is $2^{n-1}=2q$, so $b$ is well
defined on $\mathbb{Z}/2q\mathbb{Z}$ and the window is exactly one period.

\begin{lemma}[the flip]\label{lem:flip}
$b(k+q)=1-b(k)$ for every $k$.
\end{lemma}

\begin{proof}
Bit $n$ is determined modulo $2^{n+1}$. By Lemma~\ref{lem:lte},
$3^{k+q}\equiv3^{k}(1+2^{n})=3^{k}+2^{n}3^{k}\equiv3^{k}+2^{n}\pmod{2^{n+1}}$, the last
step because $3^{k}$ is odd. Adding $2^{n}$ modulo $2^{n+1}$ complements bit $n$ and leaves
the lower bits unchanged.
\end{proof}

\section{The theorem}

\begin{theorem}\label{thm:main}
Let $n\ge3$, $q=2^{n-2}$, and let $B(x)=\sum_{k=0}^{q-1}b(k)x^{k}$ be the first half of the
window. Then
\[
  A(x)=(x+1)^{q-1}\Bigl(B(x)(x+1)+x^{q}\Bigr),
\]
and the second factor is not divisible by $x+1$. Hence
$v\bigl(a(n)\bigr)=q-1=2^{n-2}-1=\mathrm{A000225}(n-2)$ exactly, and in particular
$(x+1)^{\mathrm{A000225}(n-2)}$ divides $a(n)$.
\end{theorem}

\begin{proof}
Put $S(x)=1+x+\dots+x^{q-1}$. By Lemma~\ref{lem:flip} the window of length $2q$ consists of
the block $B$ followed by its complement, and complementing a $0/1$ block of length $q$ is
adding $S$ over $\mathrm{GF}(2)$. Hence
\[
  A=B+x^{q}\bigl(B+S\bigr)=B\bigl(1+x^{q}\bigr)+S\,x^{q}.
\]
Since $q$ is a power of 2 we have $1+x^{q}=(1+x)^{q}$, and therefore
$S=(1+x^{q})/(1+x)=(1+x)^{q-1}$. Substituting,
\[
  A=B(1+x)^{q}+(1+x)^{q-1}x^{q}=(1+x)^{q-1}\Bigl(B(1+x)+x^{q}\Bigr).
\]
Evaluating the cofactor at $x=1$ over $\mathrm{GF}(2)$ gives $B(1)\cdot(1+1)+1^{q}=0+1=1
\ne0$, so $x+1$ does not divide it, and $v(A)=q-1$.
\end{proof}

\begin{proposition}[the pre-correction exponent was too large]\label{prop:old}
The statement obtained by replacing $\mathrm{A000225}(n-2)$ with
$\mathrm{A000225}(n-1)=2^{n-1}-1$ is false for every $n\ge3$. The smallest instance is
$n=3$: there $a(3)=12$, whose $\mathrm{GF}(2)$ polynomial is $x^{3}+x^{2}=x^{2}(x+1)$, so
$v(a(3))=1$, whereas $2^{n-1}-1=3$.
\end{proposition}

\begin{proof}
By Theorem~\ref{thm:main}, $v(a(n))=2^{n-2}-1<2^{n-1}-1$ for every $n\ge3$.
\end{proof}

\begin{remark}[the doubled window, and the crossrefs on the entry]
A037096 takes the same column over a window of length $2^{n}$, that is two periods rather
than one, and the entry records the crossrefs $a(n)=\lfloor
\mathrm{A037096}(n)/2^{2^{n-1}}\rfloor$ and $\mathrm{A037096}(n)=\bigl(2^{2^{n-1}}+1\bigr)
a(n)$. Both follow from Lemma~\ref{lem:flip} without any further work: over two periods the
window reads $B,\overline{B},B,\overline{B}$, so its polynomial is
$\bigl(1+x^{2q}\bigr)A(x)$, and $1+x^{2q}=x^{2^{n-1}}+1$ encodes as $2^{2^{n-1}}+1$.
Dividing by that factor recovers Theorem~\ref{thm:main}. The verification script confirms
both identities numerically as well.
\end{remark}

\begin{remark}[honest assessment]
Five lines once the flip lemma is in place, and the flip lemma is lifting the exponent. The
natural destination is an OEIS comment. Note the peculiar status of this target: the
conjecture in its currently posted form is one month old, having been corrected in July
2026, though the underlying claim dates to 1999. Proposition~\ref{prop:old} confirms that
correction was necessary and gives the smallest witness.
\end{remark}

\section{Verification}

A verification script, written separately from this note, checks every claim and prints
every number quoted. It (i) recomputes $a(2),\dots,a(9)$ from the defining sum and matches
the published DATA exactly, including $a(3)=12$, $a(4)=120$, $a(5)=57120$; (ii) verifies the
factorization of Theorem~\ref{thm:main} by explicit polynomial multiplication for
$3\le n\le15$; (iii) computes $v(a(n))$ for $3\le n\le15$, obtaining
$1,3,7,15,31,63,127,255,511,1023,2047,4095,8191$, which is $2^{n-2}-1$ in every case, with
cofactor value 1 at $x=1$ throughout, confirming exactness; and (iv) confirms the two
crossref identities $a(n)=\lfloor \mathrm{A037096}(n)/2^{2^{n-1}}\rfloor$ and
$\mathrm{A037096}(n)=(2^{2^{n-1}}+1)\,a(n)$ for $3\le n\le12$.

\begin{thebibliography}{9}
\bibitem{oeis} N.~J.~A. Sloane et al., \emph{The On-Line Encyclopedia of Integer
Sequences}, \url{https://oeis.org/A037097}, entry A037097, consulted 24 August 2026; also
A037096, A000225 and A051179.
\bibitem{rib} P.~Ribenboim, \emph{The New Book of Prime Number Records}, 3rd ed.,
Springer, 1996 (for the lifting-the-exponent lemma at $p=2$).
\end{thebibliography}
\end{document}
